\section{Fase 2: especificação da avaliação}
\label{sec:fase2}

A Fase 2 seleciona as métricas e estabelece os níveis de pontuação e os critérios de julgamento \cite{iso25040,ramos2026processo}, pelo método GQM \cite{basili1994}. Cada métrica tem definição, cálculo, fonte dos dados, procedimento de coleta e interpretação documentados na GitPage; esta seção apresenta a síntese.

\subsection{Rastreabilidade com a Fase 1}

Todos os elementos do GQM derivam de decisões da Fase 1: as quatro características priorizadas originam os objetivos G1 a G4; as subcaracterísticas mantidas originam as questões; o ponto de vista e os perfis definem os recortes de análise (por perfil em M5, M6, M7 e M11; celular em Q12; questionário por semestre); as observações preliminares fundamentam as hipóteses (OP-01 em M11, OP-02 e OP-03 em M14, OP-04 em M4, OP-05 em M6 e M7, OP-06 em M8 e M20, OP-07 em M1 e M3); as decisões apoiadas definem a forma dos resultados (ranking das tarefas, G3 para a interface responsiva, M16 a M19 por página); e o rigor intermediário define dois avaliadores em M1, M2, M10 e M20 e cinco execuções em M16 a M19.

\subsection{Objetivos de medição}

Cada objetivo segue o \textit{template} GQM (Tabela~\ref{tab:objetivos}). Em todos, o propósito é \textbf{avaliar, a fim de identificar pontos fracos}, do ponto de vista do \textbf{estudante de graduação da UnB}, com o recorte indicado. Os resultados serão apresentados na ordem de prioridade da Fase 1: G2, G1, G4 e G3.

\begin{table*}[!t]
	\caption{Objetivos de medição no template GQM}
	\label{tab:objetivos}
	\centering
	\footnotesize
	\begin{tabularx}{\textwidth}{L{2.1cm} Y Y Y Y}
		\toprule
		& \textbf{G1 -- Adequação Funcional} & \textbf{G2 -- Capacidade de Interação} & \textbf{G3 -- Flexibilidade} & \textbf{G4 -- Eficiência de Desempenho} \\
		\midrule
		\textbf{Analisar} & o portal do discente do SIGAA/UnB & a interface web do portal do discente & a interface web do portal do discente & as páginas das tarefas T1 a T6 \\
		\textbf{Com respeito à} & completude, correção e pertinência funcional & reconhecimento de adequação, apreensibilidade, operabilidade, proteção contra erros, inclusividade, autodescritividade e satisfação & adaptabilidade a navegadores, telas e toque & comportamento temporal e utilização de recursos no cliente \\
		\textbf{Recorte do ponto de vista} & estudante de graduação & calouros e pessoas com deficiência & quem acessa só pelo celular & quem acessa por rede móvel \\
		\textbf{Contexto} & T1 a T6 e planejamento do fluxo; v4.15.10.115, 2026/2 & execução de T1 a T6 em computador & Chrome, Firefox, Edge e Safari; computador, tablet e celular & medição no cliente, em rede comum e móvel simulada \\
		\bottomrule
	\end{tabularx}
	\fonte{elaborado pela equipe.}
\end{table*}

\subsection{Questões e hipóteses}

Cada questão caracteriza um atributo do foco de qualidade e recebe uma hipótese com critério de confirmação numérico, declarado antes da coleta (Tabela~\ref{tab:questoes}). Na Fase 4, cada hipótese será classificada como confirmada, refutada ou inconclusiva; uma hipótese refutada indica que a percepção da Fase 1 não se confirmou.

\begin{table*}[!t]
	\caption{Questões, atributos e hipóteses}
	\label{tab:questoes}
	\centering
	\footnotesize
	\begin{tabularx}{\textwidth}{L{0.5cm} L{5.4cm} L{2.9cm} Y}
		\toprule
		\textbf{ID} & \textbf{Questão} & \textbf{Atributo} & \textbf{Hipótese (critério de confirmação)} \\
		\midrule
		Q1 & O SIGAA oferece as funções de que o estudante precisa? & Completude & H1: cobertura entre 75\% e 90\%, pior em T4 e no planejamento do fluxo (M1 Regular) \\
		Q2 & As informações exibidas estão corretas? & Correção & H2: nenhuma divergência com as fontes oficiais (M2 = 100\%) \\
		Q3 & As funções bastam para planejar o fluxo sem ferramentas externas? & Pertinência & H3: 25\% a 50\% dos respondentes usam ferramenta externa (M3 Regular) \\
		\midrule
		Q4 & O estudante identifica onde está a função? & Reconhecimento de adequação & H4: 50\% a 75\% de acerto na primeira escolha de menu (M4 Regular) \\
		Q5 & Conclui as tarefas sem ajuda, inclusive no primeiro contato? & Apreensibilidade & H5: conclusão geral de 78\% a 90\% (M5 Bom), mas abaixo de 65\% entre calouros \\
		Q6 & A navegação evita esforço desnecessário? & Operabilidade & H6: tempo relativo de 2,5 a 4,0 e desvio de caminho de 1,33 a 2,00 (M6 e M7 Regular) \\
		Q7 & O sistema evita perda de trabalho após erros ou expiração da sessão? & Proteção contra erros & H7: até 2 de 6 cenários tratados (M20 Insatisfatório) e 50\% a 70\% dos erros recuperados (M8 Regular) \\
		Q8 & As páginas podem ser usadas por pessoas com deficiência? & Inclusividade & H8: mais de 5 erros automáticos por página (M9 Insatisfatório) e 18 a 21 de 24 verificações manuais (M10 Regular) \\
		Q9 & Compreende as informações sem legenda? & Autodescritividade & H9: 50\% a 75\% de acertos, com o código de horário de T2 abaixo de 50\% (M11 Regular) \\
		Q10 & O estudante fica satisfeito? & Satisfação (em uso) & H10: SUS médio entre 51 e 68 (M12 Regular) \\
		\midrule
		Q11 & As tarefas funcionam nos navegadores mais usados? & Adaptabilidade (software) & H11: 22 ou 23 das 24 combinações sem falha (M13 Bom) \\
		Q12 & O SIGAA se adapta a celular, tablet e toque? & Adaptabilidade (hardware) & H12: até 5 de 12 verificações adaptadas (M14 Insatisfatório) e tempo no celular 1,5 a 2,0 vezes o do computador (M15 Regular) \\
		\midrule
		Q13 & As páginas carregam em tempo aceitável, em rede comum e móvel? & Comportamento temporal & H13: LCP até 2,5~s em rede comum (M16 Bom) e acima de 2,5~s em rede móvel (M17) \\
		Q14 & O uso de recursos é compatível com rede móvel e plano limitado? & Utilização de recursos & H14: até 1~MB por página (M18 Bom), mas mais de 50 requisições (M19 Regular) \\
		\bottomrule
	\end{tabularx}
	\fonte{elaborado pela equipe.}
\end{table*}

\subsection{Métricas e níveis de pontuação}

As métricas são simples (contagens, razões ou medianas), objetivas (dois avaliadores obtêm o mesmo valor) e válidas (medem o atributo perguntado) e seguem, sempre que possível, a ISO/IEC 25023 e a ISO/IEC 25022 \cite{iso25023,iso25022}. Antes da coleta, a equipe fixa o caminho mais curto de cada tarefa e as definições de passo (clique ou toque que muda a interface), ajuda (intervenção do moderador, de colega ou de fonte externa) e tarefa concluída (resultado correto, sem ajuda, em até 5 minutos). Toda medição registra data, horário, dispositivo, navegador e rede.

Todas as métricas usam uma escala de quatro níveis: Excelente (3), Bom (2), Regular (1) e Insatisfatório (0); um valor no limite entre dois níveis pertence ao superior. Quando há referência externa, ela define os limites (Seção~\ref{sec:medidas-referencia}); nos demais casos, os limites foram definidos e justificados pela equipe. A Tabela~\ref{tab:metricas} resume as 20 métricas com seus níveis, fixados antes da coleta.

\begin{table*}[!t]
	\caption{Métricas da avaliação e níveis de pontuação}
	\label{tab:metricas}
	\centering
	\scriptsize
	\setlength{\tabcolsep}{3pt}
	\begin{tabularx}{\textwidth}{L{0.55cm} L{2.6cm} L{2.7cm} Y C{1.15cm} C{1.35cm} C{1.35cm} C{1.6cm}}
		\toprule
		\textbf{ID} & \textbf{Métrica (questão)} & \textbf{Cálculo} & \textbf{Fonte e coleta} & \textbf{Excelente} & \textbf{Bom} & \textbf{Regular} & \textbf{Insatisf.} \\
		\midrule
		M1 & Cobertura funcional (Q1) & $1-A/B$; $A$ = funções ausentes & Lista de referência (questionário, ferramentas, tarefas); 2 avaliadores & 100\% & 90--$<$100\% & 75--$<$90\% & $<$75\% \\
		M2 & Correção funcional (Q2) & $1-A/B$; $A$ = verificações divergentes & Projeto pedagógico, calendário e histórico; 2 avaliadores & 100\% & 97--$<$100\% & 90--$<$97\% & $<$90\% ou item crítico \\
		M3 & Uso de ferramentas externas (Q3) & $A/B$; $A$ = usam ferramenta & Questionário (Seção~\ref{sec:questionario}) & $\leq$10\% & $>$10--25\% & $>$25--50\% & $>$50\% \\
		M4 & Primeira escolha de menu correta (Q4) & $A/B$; $B$ = tentativas & Gravações do teste com usuários & $\geq$90\% & 75--$<$90\% & 50--$<$75\% & $<$50\% \\
		M5 & Conclusão sem ajuda (Q5) & $A/B$; $A$ = concluídas & Gravações e formulário; por perfil & $\geq$90\% & 78--$<$90\% & 60--$<$78\% & $<$60\% \\
		M6 & Tempo relativo (Q6) & $T_p/T_e$ & Gravações; especialista, 3 execuções & $\leq$1,5 & $>$1,5--2,5 & $>$2,5--4,0 & $>$4,0 \\
		M7 & Desvio de caminho (Q6) & $C_p/C_{min}$ (cliques) & Gravações; $C_{min}$ fixado antes & $\leq$1,11 & $>$1,11--1,33 & $>$1,33--2,00 & $>$2,00 \\
		M8 & Recuperação de erros (Q7) & $A/B$; $A$ = sem ajuda e sem retrabalho & Gravações do teste com usuários & $\geq$90\% & 70--$<$90\% & 50--$<$70\% & $<$50\% \\
		M20 & Cenários de erro tratados (Q7) & $A/6$ & Sessão expirada, Voltar, recarregar, nova aba, fora do período, PDF bloqueado & 6 & 5 & 3 ou 4 & 0 a 2 \\
		M9 & Erros automáticos de acessibilidade (Q8) & erros / páginas & Extensão WAVE, 6 páginas & 0 & $>$0--2 & $>$2--5 & $>$5 \\
		M10 & Critérios manuais de acessibilidade (Q8) & $A/24$ & Teclado, foco, zoom 200\%, leitor de tela; 2 avaliadores & 24 & 22 ou 23 & 18 a 21 & $\leq$17 \\
		M11 & Interpretação sem legenda (Q9) & $A/B$; 3 perguntas & Perguntas em T2, T4 e T6 & $\geq$90\% & 75--$<$90\% & 50--$<$75\% & $<$50\% \\
		M12 & Pontuação SUS (Q10) & SUS, 0 a 100 & Questionário após o teste \cite{brooke1996} & $\geq$80,3 & 68--$<$80,3 & 51--$<$68 & $<$51 \\
		M13 & Compatibilidade de navegadores (Q11) & $A/24$ (6 tarefas $\times$ 4) & Execução pela equipe, com captura & 24 & 22 ou 23 & 18 a 21 & $\leq$17 \\
		M14 & Páginas adaptadas (Q12) & $A/12$ (6 páginas $\times$ 2) & 360 e 768~px: rolagem, fonte, sobreposição, toque, alvo de 24~px & 12 & 10 ou 11 & 6 a 9 & $\leq$5 \\
		M15 & Tempo relativo no celular (Q12) & $T_{cel}/T_{comp}$ & Execução cronometrada & $\leq$1,2 & $>$1,2--1,5 & $>$1,5--2,0 & $>$2,0 \\
		M16 & LCP em rede comum (Q13) & média das medianas (s) & Lighthouse \textit{Desktop}, 5 execuções \cite{lighthouse} & $\leq$1,0 & $>$1,0--2,5 & $>$2,5--4,0 & $>$4,0 \\
		M17 & LCP em rede móvel (Q13) & média das medianas (s) & Lighthouse \textit{Mobile} & $\leq$1,0 & $>$1,0--2,5 & $>$2,5--4,0 & $>$4,0 \\
		M18 & Peso por página (Q14) & média das medianas (MB) & Lighthouse (\textit{total byte weight}) & $\leq$0,5 & $>$0,5--1,0 & $>$1,0--2,0 & $>$2,0 \\
		M19 & Requisições por página (Q14) & média das medianas & Lighthouse (\textit{network requests}) & $\leq$25 & $>$25--50 & $>$50--80 & $>$80 \\
		\bottomrule
	\end{tabularx}
	\fonte{elaborado pela equipe, com base em \cite{iso25023,iso25022,sauro2011,sauro2016,wcag22,webvitals,nielsen1993}.}
\end{table*}

\subsection{Critérios de julgamento}

O julgamento vai do nível quantitativo ao conceitual. (1) A \textbf{questão} recebe o \textbf{menor} nível entre as suas métricas: um atributo só é Bom se todas as medidas forem ao menos Boas. (2) O \textbf{objetivo} recebe a \textbf{média} dos níveis das questões, para que subcaracterísticas com mais métricas não pesem mais. (3) A \textbf{característica} \textit{atende plenamente} se a média for $\geq 2{,}5$ sem questão abaixo de Bom; \textit{atende} se for $\geq 2{,}0$ sem questão Insatisfatória; \textit{atende parcialmente} se for $\geq 1{,}0$; \textit{não atende} se for menor; e é \textit{inconclusiva} se mais da metade das questões ficar sem dados.

Regras complementares: métricas sem dados ficam fora do cálculo; divergência em menção, IRA ou carga horária pendente leva M2 a Insatisfatório; recortes por perfil, página ou tarefa localizam problemas, mas não alteram o nível, e recortes com menos de três participantes são descritivos; cada tarefa recebe um índice (média dos níveis de M5, M6, M7 e M1 na tarefa) para ordenar as tarefas da pior para a melhor; e \textbf{não é calculada nota única} para o SIGAA, pois uma média geral esconderia os problemas a priorizar.

\subsection{Hierarquia GQM}

A Fig.~\ref{fig:gqm} apresenta a hierarquia completa nos três níveis do GQM. Toda métrica está ligada a uma questão, e toda questão, a um objetivo e a uma hipótese.

\begin{figure*}[!t]
	\centering
	\begin{tikzpicture}[
		y=-0.5cm,
		g/.style={draw=blue!60!black, fill=blue!12, rounded corners, align=center, text width=3.0cm, font=\scriptsize, inner sep=2pt},
		q/.style={draw=teal!60!black, fill=teal!8, rounded corners, align=left, text width=5.6cm, font=\scriptsize, inner sep=1.5pt},
		m/.style={draw=orange!70!black, fill=orange!8, rounded corners, align=left, text width=5.8cm, font=\scriptsize, inner sep=1.5pt},
		lig/.style={-{Stealth[length=1.4mm]}, gray!70!black}]
		\node[font=\footnotesize\bfseries] at (0,-1.3) {Nível conceitual (objetivos)};
		\node[font=\footnotesize\bfseries] at (5.6,-1.3) {Nível operacional (questões)};
		\node[font=\footnotesize\bfseries] at (12.4,-1.3) {Nível quantitativo (métricas)};
		\foreach \id/\txt/\y in {
			M1/Cobertura funcional/0, M2/Correção funcional/1, M3/Uso de ferramentas externas/2,
			M4/Acerto na primeira escolha de menu/3, M5/Conclusão sem ajuda/4, M6/Tempo relativo de execução/5,
			M7/Desvio de caminho/6, M8/Recuperação de erros/7, M20/Cenários de erro tratados/8,
			M9/Erros automáticos de acessibilidade/9, M10/Critérios manuais de acessibilidade/10,
			M11/Interpretação sem legenda/11, M12/Pontuação SUS/12, M13/Compatibilidade de navegadores/13,
			M14/Páginas adaptadas a telas pequenas e toque/14, M15/Tempo relativo no celular/15, M16/LCP em rede comum/16,
			M17/LCP em rede móvel simulada/17, M18/Peso transferido por página/18, M19/Requisições por página/19}
			\node[m] (\id) at (12.4,\y) {\textbf{\id} -- \txt};
		\foreach \id/\txt/\y in {
			Q1/Completude funcional/0, Q2/Correção funcional/1, Q3/Pertinência funcional/2,
			Q4/Reconhecimento de adequação/3, Q5/Apreensibilidade/4, Q6/Operabilidade/5.5,
			Q7/Proteção contra erros do usuário/7.5, Q8/Inclusividade/9.5, Q9/Autodescritividade/11,
			Q10/Satisfação (qualidade em uso)/12, Q11/Adaptabilidade a navegadores/13,
			Q12/Adaptabilidade a telas e toque/14.5, Q13/Comportamento temporal/16.5,
			Q14/Utilização de recursos/18.5}
			\node[q] (\id) at (5.6,\y) {\textbf{\id} -- \txt};
		\node[g] (G1) at (0,1)     {\textbf{G1} Adequação Funcional};
		\node[g] (G2) at (0,7.5)   {\textbf{G2} Capacidade de Interação (Usabilidade)};
		\node[g] (G3) at (0,13.75) {\textbf{G3} Flexibilidade (Portabilidade)};
		\node[g] (G4) at (0,17.5)  {\textbf{G4} Eficiência de Desempenho};
		\foreach \gg/\qq in {G1/Q1, G1/Q2, G1/Q3, G2/Q4, G2/Q5, G2/Q6, G2/Q7, G2/Q8, G2/Q9, G2/Q10,
			G3/Q11, G3/Q12, G4/Q13, G4/Q14}
			\draw[lig] (\gg.east) -- (\qq.west);
		\foreach \qq/\mm in {Q1/M1, Q2/M2, Q3/M3, Q4/M4, Q5/M5, Q6/M6, Q6/M7, Q7/M8, Q7/M20, Q8/M9,
			Q8/M10, Q9/M11, Q10/M12, Q11/M13, Q12/M14, Q12/M15, Q13/M16, Q13/M17, Q14/M18, Q14/M19}
			\draw[lig] (\qq.east) -- (\mm.west);
	\end{tikzpicture}
	\caption{Hierarquia GQM da avaliação. Fonte: elaborado pela equipe.}
	\label{fig:gqm}
\end{figure*}

\subsection{Relações entre as características}

O roteiro exige discutir, ao final, as relações entre as características \cite{ramos2026roteiro}. As relações esperadas na Fase 1 viraram quatro hipóteses, sem nível de pontuação, que serão declaradas estabelecidas ou não na Fase 4 (Tabela~\ref{tab:relacoes}).

\begin{table}[!t]
	\caption{Hipóteses de relação entre as características}
	\label{tab:relacoes}
	\centering
	\footnotesize
	\begin{tabularx}{\columnwidth}{L{0.6cm} L{2.0cm} Y}
		\toprule
		\textbf{ID} & \textbf{Relação} & \textbf{Hipótese e verificação} \\
		\midrule
		HR1 & Flexibilidade $\rightarrow$ Cap. Interação & O esforço extra no celular se concentra nas páginas não adaptadas: M15 por tarefa cruzado com M14 \\
		HR2 & Desempenho $\rightarrow$ Cap. Interação & Tarefas com maior LCP têm maior tempo relativo: ordenações de M16 e M6 (Spearman descritivo) \\
		HR3 & Cap. Interação $\rightarrow$ Adequação Funcional & Funções existentes são declaradas inexistentes: métrica auxiliar MA1 \\
		HR4 & Desempenho $\rightarrow$ Flexibilidade & O LCP em rede móvel é ao menos o dobro do LCP em rede comum: razão M17/M16 \\
		\bottomrule
	\end{tabularx}
	\fonte{elaborado pela equipe.}
\end{table}

\subsection{Questionário com estudantes de diferentes semestres}
\label{sec:questionario}

O teste com usuários mede o comportamento em profundidade, mas com poucos participantes. Para ampliar a amostra e comparar estudantes em diferentes momentos do curso, a especificação inclui um questionário online anônimo, de 5 a 8 minutos, com seis blocos: perfil (semestre, campus, previsão de formatura), acesso (dispositivo, rede e problemas no celular), facilidade percebida em T1 a T6 (escala SEQ de 1 a 5 \cite{sauro2016}), compreensão sem legenda (código de horário ``35T23'' e sigla IRA), satisfação (SUS \cite{brooke1996}) e ferramentas externas. Ele é a fonte de M3 e de parte da lista de referência de M1 e produz indicadores auxiliares de percepção (IQ1 a IQ6) que, como a MA1, não recebem nível nem alteram o julgamento. Esses indicadores verificam a hipótese de perfil \textbf{HP1}: estudantes de 1º e 2º semestres relatam menor facilidade e acertam menos o código de horário do que estudantes a partir do 5º semestre. A comparação usa medianas e proporções e, com ao menos dez respondentes por grupo, os testes de Mann-Whitney e exato de Fisher, com significância de 5\%; caso contrário, é descritiva. O instrumento completo está na GitPage.
